\usepackage{ctex}                                   % 中文支持
\usepackage{amsmath, amssymb, amsthm, mathtools}    % 数学环境
\usepackage{mathrsfs}                               % 数学花体字支持
\usepackage{makeidx}                                % 索引支持
\usepackage{enumitem}                               % 列表环境
\usepackage{tikz}                                   % 绘图支持
% \usepackage{pgfplots}                               % 绘图支持
% \usepackage{pgfplotstable}                          % 表格支持
\usepackage{bm}                                     % 粗体数学符号
\usepackage{geometry}                               % 页面设置
\usepackage{xcolor}                                 % 颜色支持
\usepackage{graphicx}                               % 图片支持
\usepackage{tcolorbox}                              % 盒子支持

% 参考文献支持
\usepackage[
  backend=biber,
  style=numeric,
  citestyle=numeric,
  sorting=none,
  maxnames=3,
  minnames=1,
  maxbibnames=99,
  minbibnames=1,
  giveninits=true,
  doi=false,
  isbn=false,
  url=true
]{biblatex}

% 超链接支持
\usepackage[
    colorlinks=true,        % 启用彩色链接
    linkcolor=blue,         % 章节、图表等内部链接颜色
    citecolor=blue,         % 引用链接颜色
    urlcolor=blue           % URL 链接颜色
    pdfborder={0 0 0},      % 去掉链接周围的边框
]{hyperref}

\usepackage{cleveref}                               % 智能引用支持

%--------------------------------------------------------------------------------------
% 数学环境设置
%--------------------------------------------------------------------------------------
\numberwithin{equation}{section}            % 公式编号按章节重新计数，例如 chapter2 section3 的第 7 个公式编号为 (2.3.7)

\theoremstyle{definition}                   % 定义环境样式，适用于定义、例等
\newtheorem{definition}{定义}[section]      % 定义环境，按章节重新计数，例如 chapter2 section3 的第 1 个定义编号为 (2.3.1)
\newtheorem{example}{例}[section]           % 例环境，按章节重新计数，例如 chapter2 section3 的第 1 个例编号为 (2.3.1)

\theoremstyle{plain}                        % 定理环境样式，plain 是默认样式，适用于定理、命题、引理、推论等
\newtheorem{axiom}{公理}[section]           % 公理环境，按章节重新计数，例如 chapter2 section3 的第 1 个公理编号为 (2.3.1)
\newtheorem{theorem}{定理}[section]         % 定理环境，按章节重新计数，例如 chapter2 section3 的第 1 个定理编号为 (2.3.1)
\newtheorem{proposition}{命题}[section]     % 命题环境，按章节重新计数，例如 chapter2 section3 的第 1 个命题编号为 (2.3.1)
\newtheorem{lemma}{引理}[section]           % 引理环境，按章节重新计数，例如 chapter2 section3 的第 1 个引理编号为 (2.3.1)
\newtheorem{corollary}{推论}[section]       % 推论环境，按章节重新计数，例如 chapter2 section3 的第 1 个推论编号为 (2.3.1)

\theoremstyle{remark}                       % 评注环境样式，适用于评注、注释等
\newtheorem{remark}{注}[section]          % 评注环境，按章节重新计数，例如 chapter2 section3 的第 1 个评注编号为 (2.3.1)
\newtheorem{note}{说明}[section]          % 说明环境，按章节重新计数，例如 chapter2 section3 的第 1 个说明编号为 (2.3.1)

% 设置 \newtheorem 块的前后间隔
\AtBeginEnvironment{axiom}{\vspace{0.8em}}
\AtEndEnvironment{axiom}{\vspace{0.0em}}
\AtBeginEnvironment{definition}{\vspace{0.8em}}
\AtEndEnvironment{definition}{\vspace{0.0em}}
\AtBeginEnvironment{theorem}{\vspace{0.8em}}
\AtEndEnvironment{theorem}{\vspace{0.0em}}
\AtBeginEnvironment{proposition}{\vspace{0.8em}}
\AtEndEnvironment{proposition}{\vspace{0.0em}}
\AtBeginEnvironment{lemma}{\vspace{0.8em}}
\AtEndEnvironment{lemma}{\vspace{0.0em}}
\AtBeginEnvironment{corollary}{\vspace{0.8em}}
\AtEndEnvironment{corollary}{\vspace{0.0em}}
\AtBeginEnvironment{example}{\vspace{0.8em}}
\AtEndEnvironment{example}{\vspace{0.0em}}
\AtBeginEnvironment{remark}{\vspace{0.8em}}
\AtEndEnvironment{remark}{\vspace{0.0em}}
\AtBeginEnvironment{note}{\vspace{0.8em}}
\AtEndEnvironment{note}{\vspace{0.0em}}

% 设置 proof 块的前后间距
\AtBeginEnvironment{proof}{\vspace{0.0em}}
\AtEndEnvironment{proof}{\vspace{0.0em}}

%--------------------------------------------------------------------------------------
% cref 链接环境设置
%--------------------------------------------------------------------------------------
\crefname{axiom}{公理}{公理}               % cref 引用 axiom 环境时使用 "公理" 作为单数和复数形式
\crefname{definition}{定义}{定义}         % cref 引用 definition 环境时使用 "定义" 作为单数和复数形式
\crefname{theorem}{定理}{定理}           % cref 引用 theorem 环境时使用 "定理" 作为单数和复数形式
\crefname{proposition}{命题}{命题}       % cref 引用 proposition 环境时使用 "命题" 作为单数和复数形式
\crefname{lemma}{引理}{引理}             % cref 引用 lemma 环境时使用 "引理" 作为单数和复数形式
\crefname{corollary}{推论}{推论}       % cref 引用 corollary 环境时使用 "推论" 作为单数和复数形式
\crefname{example}{例}{例}             % cref 引用 example 环境时使用 "例" 作为单数和复数形式
\crefname{remark}{注}{注}             % cref 引用 remark 环境时使用 "注" 作为单数和复数形式
\crefname{note}{说明}{说明}             % cref 引用 note 环境时使用 "说明" 作为单数和复数形式
\crefname{section}{节}{节}             % cref 引用 section 环境时使用 "节" 作为单数和复数形式
\crefname{chapter}{章}{章}             % cref 引用 chapter 环境时使用 "章" 作为单数和复数形式
\crefname{equation}{式}{式}           % cref 引用 equation 环境时使用 "式" 作为单数和复数形式
\crefname{figure}{图}{图}             % cref 引用 figure 环境时使用 "图" 作为单数和复数形式
\crefname{table}{表}{表}             % cref 引用 table 环境时使用 "表" 作为单数和复数形式

%--------------------------------------------------------------------------------------
% 页面布局设置
%--------------------------------------------------------------------------------------
\geometry{
    left=1.5cm,    % 左边距
    right=1.5cm,   % 右边距
    top=2.5cm,     % 上边距
    bottom=3cm     % 下边距
}

\setlength{\parskip}{0em} % 段间距
\setlength{\parindent}{0em} % 首行缩进

%--------------------------------------------------------------------------------------
% 枚举块设置
%--------------------------------------------------------------------------------------
\setlist[enumerate]
{
    topsep=0.3em, % 块前后间隔
    itemsep=0.3em, % 枚举项之间的间隔
    partopsep=0em, % 段落前间隔
    parsep=0em    % 段落内间隔
}

%--------------------------------------------------------------------------------------
% 参考文献设置
%--------------------------------------------------------------------------------------
\addbibresource{bibliography/references.bib} % 添加参考文献数据库文件

%--------------------------------------------------------------------------------------
% 注意盒子
%--------------------------------------------------------------------------------------
\newtcolorbox{important}[1][]
{
    colback=gray!10!white, % 盒子背景色
    colframe=gray!50!black, % 盒子边框色
    title=说明, % 盒子标题
    before=\vspace{0.5em}, % 盒子前间距
    after=\vspace{1em}, % 盒子后间距
    #1
}